\section{Derived induction and the trace obstruction}
\label{b:section}

We now construct the bimodule that will join the two twists in the
triangular algebra.  The two-factor filtration follows the organization
of \cite[Section~6]{OpenAIAR2026}. A related use of tensor products of
bimodules attached to Hochschild classes occurs in support-variety
theory~\cite[Definition~2.1 and Theorem~2.3]{ErdmannEtAl2003}.
Here the bimodule comes from derived
induction, and its evaluated lifting property is proved by a trace
obstruction.

\begin{proposition}\label{b:bimodule}
There is a finite \(E\)-bimodule \(Y\), projective separately as a left and
as a right \(E\)-module, such that
\begin{equation}\label{b:W}
 W^a:=\sHom_E(X,(YX)[a])
 \cong
 \begin{cases}
 k,&a=-3,0,\\
 0,&a\notin\{-3,0\}.
 \end{cases}
\end{equation}
For every \(H\in k^\times\), there is an ordinary bimodule map
\(g_H:U_H\to Y\) whose evaluation, under \(U_HX=X\), represents a nonzero
element of \(W^0\).  The same \(Y\) works for all \(H\), and all modules and
maps are defined over \(k\).
\end{proposition}

The two assertions require different arguments. A filtration of derived
induction will compute the groups \(W^a\). In particular, \(W^0=k\)
provides nonzero stable \(E\)-module maps \(X\to YX\), but does not show
that one is the evaluation of a bimodule map \(U_H\to Y\). The trace
obstruction, combined with stable duality, will prove that this evaluation
has nonzero image for every \(H\).

\subsection{A finite representative of derived induction}

Let \(K\) be the two-sided bar complex for derived induction:
\[
 K_i=E\otimes B^{\otimes i}\otimes E\quad(i\geq0),\qquad K_i=0\quad(i<0).
\]
This is the ordinary-algebra instance of the bar model for derived
tensor functors; see~\cite[Sections~6.1 and~6.6]{Keller1994}.
The finite representative needed here will follow from the separate
projective-dimension bounds on the two sides.
For \(i\geq1\), its differential is
\begin{align*}
 d_i(a\otimes b_1\otimes\cdots\otimes b_i\otimes c)
 ={}&
 ab_1\otimes b_2\otimes\cdots\otimes b_i\otimes c\\
 &+\sum_{r=1}^{i-1}
 a\otimes b_1\otimes\cdots\otimes b_rb_{r+1}
       \otimes\cdots\otimes b_i\otimes c\\
 &+a\otimes b_1\otimes\cdots\otimes b_{i-1}\otimes b_ic .
\end{align*}
As throughout, signs disappear in characteristic two.  Associativity
gives \(d_{i-1}d_i=0\), since each double multiplication occurs twice.
Multiplication in degree zero defines a chain map \(m:K\to E\).

\begin{lemma}\label{b:bar-model}
The complex \(K\) represents \(E\otimes_B^{\mathbf L}E\), has finite
projective \(E^e\)-terms and has homology only in degrees \(0,\ldots,4\).
The bimodule \(c_5K\) is projective on both sides.  Consequently its
stable image
\[
 \cI=\langle K\rangle=(c_5K)[5]
\]
has finite side-projective representatives, as does \(\cI[-1]\).
One such representative \(Y\) of \(\cI[-1]\) is obtained by taking four
successive finite bimodule cosyzygies of \(c_5K\).
\end{lemma}

\begin{proof}
The usual bar resolution of \(B\) has terms
\(B\otimes B^{\otimes i}\otimes B\).  Its augmented complex is contracted
as a complex of vector spaces by inserting a unit.  Inserting a unit at
one end also proves that tensoring it with a module at the other end
gives a projective resolution of that module on the remaining side.
Inducing the two outer actions gives \(K\), with each term a finite free
\(E^e\)-module.

On the left, \(B\otimes B^{\otimes i}\otimes E\) is a projective
\(B\)-resolution of \({}_BE\).  By Proposition~\ref{a:dual-dimensions}, this
module has a projective resolution supported in degrees \(0,\ldots,4\).
Comparison of projective resolutions gives a homotopy equivalence
between these two resolutions.  Applying \(E\otimes_B-\) preserves the
homotopy equivalence, so \(K\), as a left \(E\)-complex, is homotopy
equivalent to a complex of projectives supported in those degrees.
The same argument using \(E_B\) works on the right.  In particular
\(H_i(K)=0\) for \(i>4\).

For completeness, this equivalence proves projectivity of the specific
cokernel \(c_5K\).  On either side, choose a homotopy from the identity
of \(K\) to a map factoring through the complex ending in degree \(4\).
In degree \(5\) it gives
\[
 1_{K_5}=d_6h_5+h_4d_5.
\]
Modulo \(\im d_6\), the identity of \(c_5K\) therefore factors as
\[
 c_5K\xrightarrow{\ \overline d_5\ }K_4
       \xrightarrow{\ q_5h_4\ }c_5K,
\]
where \(q_5:K_5\to c_5K\) is the quotient map.  Thus \(c_5K\) is a
summand of a projective on each side.  This assertion concerns the two
side-module structures; the two splittings need not be bimodule
splittings.  The high exact sequences
\(0\to c_{n+1}K\to K_n\to c_nK\to0\), for \(n\geq5\), now split
on each side.  Induction shows that all higher cokernels are
side-projective as well.

To make the finite choice explicit, put \(N_0=c_5K\), and for
\(r=0,1,2,3\) use
\[
 0\longrightarrow N_r\longrightarrow
 \Hom_k(E^e,N_r)\longrightarrow N_{r+1}\longrightarrow0,\qquad
 n\longmapsto(a\longmapsto an).
\]
Here the \(E^e\)-action on the middle term is by right multiplication
on its input.  The middle term is finite injective, hence projective
because \(E^e\) is symmetric.  Since \(N_r\) is projective and therefore
injective on each \(E\)-side, this sequence splits on both sides; it
follows inductively that every \(N_r\) is side-projective.
Take \(Y=N_4\).  Its stable class is
\((c_5K)[4]=\cI[-1]\), as required.  Lemmas~\ref{h:stable-image}
and~\ref{h:evaluation} give the stated stable images and their
evaluation compatibilities.  We denote the stable map induced by
\(m\) by \(\overline m:\cI\to E\).
\end{proof}

\subsection{The stable Hom profile in every degree}

\begin{lemma}\label{b:profile}
The bimodule \(Y\) of Lemma~\ref{b:bar-model} satisfies
\eqref{b:W}.
\end{lemma}

\begin{proof}
The evaluated complex is
\[
 K\otimes_E X=E\otimes B^{\otimes\bullet}\otimes X.
\]
Before inducing to \(E\), its bar model and \(R\otimes R\) are
projective \(B\)-resolutions of \(X=s\otimes s\).  They are homotopy
equivalent.  Induction preserves this equivalence and gives
\begin{equation}\label{b:evaluated-model}
 K\otimes_E X\simeq
 (T\otimes_C R)\otimes(T\otimes_C R)
 =\mathcal J\otimes\mathcal J .
\end{equation}
In particular, Lemma~\ref{h:evaluation} identifies its stable image
with \(\cI X\).

We retain the actual attaching maps in the triangle
\[
 s[2]\longrightarrow\mathcal J\longrightarrow s
       \xrightarrow{\tau}s[3]
\]
of \eqref{e:triangle}.  Choose bounded-below projective
models \(A_\bullet\) for \(s[2]\) and \(Q_\bullet\) for \(s\), and a
chain map \(t:Q_\bullet\to A_\bullet[1]\) representing \(\tau\).
The cone of \(t[-1]\) has underlying graded complex
\(A_\bullet\oplus Q_\bullet\), differential
\[
 \begin{pmatrix}d_A&t\\0&d_Q\end{pmatrix},
\]
and represents \(\mathcal J\), by Lemma~\ref{h:derived-models}.
Indeed the cone construction for the rotated triangle gives this
identification; comparisons of projective
models lift the connecting map and give homotopy equivalences of the
cones.  Thus this is a homotopy equivalent projective model, with its
specified subcomplex and quotient, rather than only an equality of
homology groups.

Tensor two copies of this model and shift by \([-4]\).
Filter first by \(A_\bullet\otimes A_\bullet\), then add
\(Q_\bullet\otimes A_\bullet\) and \(A_\bullet\otimes Q_\bullet\),
and finally add \(Q_\bullet\otimes Q_\bullet\).
Every inclusion is split in each graded degree, and the successive
quotients resolve
\begin{equation}\label{b:three-levels}
 X,\qquad X[-2]\oplus X[-2],\qquad X[-4].
\end{equation}
The stable image of the full complex is
\(\mathcal L=(\cI X)[-4]\).
With the middle terms in the displayed order, the adjacent
connecting maps are the shifted row and column
\[
  (\tau_1,\tau_2),\qquad
  \binom{\tau_2}{\tau_1}.
\]
These are the cross differentials \(t\) in the relevant factor and
the identity in the other factor.  Lemma~\ref{h:stable-image}
identifies them with the connecting maps after taking stable images.

We first compute the kernels and cokernels of these adjacent maps in
every degree. On the polynomial groups the maps form the two-variable
Koszul sequence; on the negative groups they form its dual. We then use the
two long exact sequences of the filtration to control the full top
attaching map, including its possible component into the bottom level.

Write
\(\mathcal P_m=k[\tau_1,\tau_2]_m\) for polynomials of total exponent
\(m\), and put \(\mathcal P_m=0\) for \(m<0\).
By Proposition~\ref{e:all-degrees},
\[
 \cH^{3m}=\mathcal P_m,\qquad
 \cH^{-3m-1}=D\mathcal P_m\quad(m\geq0),
\]
and all other \(\cH^a\) vanish.  Multiplication by a degree-three
class on the negative groups is transpose to polynomial
multiplication.  For each integer \(l\), denote the adjacent maps by
\[
 \begin{aligned}
 d_l:(\cH^{l-2})^2&\longrightarrow\cH^{l+1},&
 d_l(u,v)&=\tau_1u+\tau_2v,\\
 d'_l:\cH^{l-4}&\longrightarrow(\cH^{l-1})^2,&
 d'_l(w)&=(\tau_2w,\tau_1w).
 \end{aligned}
\]
Products in these formulas mean shifted postcomposition.  We claim
\begin{align}
 \ker d_l&=\im d'_{l-1},\label{b:koszul-middle}\\
 \coker d_{l-1}&=
 \begin{cases}k,&l=0,\\0,&l\ne0,\end{cases}
 \qquad
 \ker d'_l=
 \begin{cases}k,&l=3,\\0,&l\ne3.\end{cases}
 \label{b:koszul-ends}
\end{align}

To check the middle identity, the polynomial calculation for \(m\geq0\)
is the exact sequence
\begin{equation}\label{b:polynomial-koszul}
 0\longrightarrow\mathcal P_{m-1}
 \xrightarrow{\,(\tau_2,\tau_1)^{\mathsf t}\,}
 \mathcal P_m^2
 \xrightarrow{\,(\tau_1,\tau_2)\,}
 \mathcal P_{m+1}\longrightarrow0 .
\end{equation}
The row is onto since every positive-degree monomial is divisible
by a variable.  If \(\tau_1u+\tau_2v=0\), coprimality implies
\(u=\tau_2w\) and \(v=\tau_1w\), uniquely; the signs agree in
characteristic two.  When the middle group is
\((D\mathcal P_m)^2\), use the dual of
\eqref{b:polynomial-koszul}, interchanging the two middle slots:
\[
 D\mathcal P_{m+1}\longrightarrow
 (D\mathcal P_m)^2\longrightarrow D\mathcal P_{m-1}.
\]
At \(m=0\), the last term is the actual group \(\cH^2=0\), and the
first map \(D\mathcal P_1\to(D\mathcal P_0)^2\) is an isomorphism.
This is the seam \(l=1\).  If the middle group is unsupported, both
its kernel and the image into it are zero.  This proves
\eqref{b:koszul-middle} for every \(l\).

For the row cokernel, a target \(\cH^l=\mathcal P_m\) gives
\(\mathcal P_{m-1}^2\to\mathcal P_m\), which is onto for \(m\geq1\)
and has cokernel \(k\) for \(m=0\).  A negative target
\(\cH^l=D\mathcal P_m\) gives the dual of the injective column
\(\mathcal P_m\to\mathcal P_{m+1}^2\), and is onto for every
\(m\geq0\).  All other targets are zero.
For the column kernel, a nonnegative source \(\mathcal P_m\) maps
injectively into \(\mathcal P_{m+1}^2\).
A source \(D\mathcal P_m\) for \(m\geq1\) is the dual of the
surjective row \(\mathcal P_{m-1}^2\to\mathcal P_m\), so again the
column is injective.  The remaining source is
\(\cH^{-1}=D\mathcal P_0=k\), at \(l=3\), with zero target
\((\cH^2)^2\).  These cases prove \eqref{b:koszul-ends}.
They include the only possible multiplication from a supported
negative degree across zero, namely \(\cH^{-1}\to\cH^2=0\).

It remains to pass from adjacent maps to the whole filtered
object.  Let \(\mathcal L_1\) be the stable image of the lower two
levels of \eqref{b:three-levels}, and put
\(N^l=\sHom_E(X,\mathcal L_1[l])\).
The first two levels give
\begin{equation}\label{b:first-two-levels}
 0\longrightarrow\coker d_{l-1}\longrightarrow N^l
 \xrightarrow{q_l}\ker d_l\longrightarrow0.
\end{equation}
Let
\[
 r_l:\cH^{l-4}\longrightarrow N^{l+1}
\]
be the \emph{actual} connecting map obtained by attaching the top
level.  Naturality of the quotient to the middle level gives
\(q_{l+1}r_l=d'_l\).
Thus \(\ker r_l\subseteq\ker d'_l\), which is zero unless \(l=3\).
At that index \(N^4=0\): its left term in
\eqref{b:first-two-levels} is \(\coker d_3=0\), and its right term
is a kernel in \((\cH^2)^2=0\).  Hence
\(\ker r_3=\cH^{-1}=k\).

Moreover,
\[
 q_l(\im r_{l-1})=\im d'_{l-1}=\ker d_l.
\]
For \(l\ne0\), \(q_l\) is an isomorphism, so \(r_{l-1}\) is onto.
At \(l=0\), the vanishing \(\cH^{-2}=0\) gives \(N^0=k\), while
the whole incoming source of \(r_{-1}\) is \(\cH^{-5}=0\).
Consequently \(\coker r_{l-1}\) is \(k\) at \(l=0\) and zero
elsewhere.

The final long exact sequence gives
\[
 0\longrightarrow\coker r_{l-1}
 \longrightarrow\sHom_E(X,\mathcal L[l])
 \longrightarrow\ker r_l\longrightarrow0.
\]
The two scalar groups therefore survive precisely at \(l=0,3\).
This argument permits an additional component of the top attaching
map into the bottom level: \(N^4=0\) and \(\cH^{-5}=0\) already
force the asserted kernel and cokernel for the full \(r_l\).
Finally \(YX=(\cI X)[-1]=\mathcal L[3]\), so \(l=a+3\).
This gives exactly \eqref{b:W}.
\end{proof}

\subsection{From perfect factorizations to an ordinary trace}

The required Hom profile is now established. It remains to construct
bimodule maps from the twists whose evaluations generate its degree-zero
line. We first show that a particular character map cannot factor
through a perfect complex; stable duality will then turn this
obstruction into the required maps.

Let \(\chi:E\to k\) be the character of \(X\), write
\(\chi_B=\chi|_B\), and retain \(p_f=f\otimes f\).
The \(B\)-bimodule structure on \(DB\) is
\((a\lambda b)(c)=\lambda(bca)\).

We use the bimodule \(DB\otimes B\), with its outer actions, to connect
perfect factorizations with ordinary traces. Using the bound
\(\idim_B B\leq4\), left coinduction will replace a perfect complex by
a bounded complex with terms in \(\add(DB\otimes B)\). For an ordinary
composite \(v:B\to DB\) through \(DB\otimes B\), the identification
\(DB\otimes B\simeq\Hom_k(B,B)\) will express the scalar \(v(1)(p_f)\)
as the trace of a linear operator.

\begin{lemma}\label{b:ordinary-factorization}
If a bimodule map \(v:B\to DB\) factors in the derived category
through a perfect \(B^e\)-complex, then it factors by ordinary
bimodule maps through an object of \(\add(DB\otimes B)\), where the
tensor product has its outer bimodule actions.
\end{lemma}

\begin{proof}
Represent the perfect intermediate object by a bounded complex
\(P\) of finite projective bimodules.  We first replace it by a
bounded complex \(P'\) whose terms belong to \(\add(DB\otimes B)\).

For a finite bimodule \(L\), define
\[
 J_B(L)=\Hom_k(B,L),\qquad
 (a\phi b)(c)=\phi(ca)b,\qquad
 j_L(z)(c)=cz.
\]
The map \(j_L:L\to J_B(L)\) is a natural bimodule monomorphism.
Set \(C^0(L)=L\) and
\(C^{r+1}(L)=\coker j_{C^r(L)}\).  These are additive functors
and preserve split retracts.  For a fixed \(d=4\) we have the
functorial exact sequence
\begin{equation}\label{b:coresolution}
 0\longrightarrow L\longrightarrow J_B(C^0(L))
 \longrightarrow\cdots\longrightarrow J_B(C^{d-1}(L))
 \longrightarrow C^d(L)\longrightarrow0.
\end{equation}

For the free bimodule \(L=B\otimes B\), this construction is the
left-module coinduced coresolution of \(B\), tensored over \(k\)
with the right regular module \(B\).  This follows from
\(\Hom_k(B,B\otimes B)=\Hom_k(B,B)\otimes B\) and from exactness
of tensoring vector spaces.  Each coinduced left-module term is
a finite injective, hence belongs to \(\add(DB)\).
The last left-module cokernel is also injective: by
Proposition~\ref{a:dual-dimensions}, \(\idim({}_BB)\leq4\), and
dimension shifting gives, for every finite left module \(V\),
\[
 \Ext_B^1(V,C^d(B))
 \cong\Ext_B^{d+1}(V,B)=0.
\]
Here \(C^d(B)\) denotes the same coresolution construction on the
left regular module.  Thus every term on the right of \(L\) in
\eqref{b:coresolution}, including the last cokernel, lies in
\(\add(DB\otimes B)\) when \(L=B\otimes B\).
The assertion extends to every finite projective bimodule
\(L\), since the functors in \eqref{b:coresolution} preserve finite
sums and split retracts.

Apply this one fixed finite coresolution to each term of \(P\).
Naturality provides the horizontal maps, giving a bounded double
complex.  Put the first coinduced term in homological degree \(0\)
relative to \(L\), and the last cokernel in degree \(-d\).
Its totalization \(P'\) is bounded and has all terms in
\(\add(DB\otimes B)\).  The augmentation \(P\to P'\) is a
quasi-isomorphism: filtering its cone by the finitely many original
degrees has exact associated complexes
\eqref{b:coresolution}.  Thus the given derived factorization can
be written \(B\to P'\to DB\).

For \(V\in\add(DB\otimes B)\) and \(i>0\), we next prove
\begin{equation}\label{b:two-orthogonalities}
 \Ext_{B^e}^i(B,V)=0,\qquad
 \Ext_{B^e}^i(V,DB)=0.
\end{equation}
It suffices to take \(V=DB\otimes B\).
Coinduction identifies, for a bimodule \(Q\),
\[
 \Hom_{B^e}(Q,DB\otimes B)
 \cong\Hom_{B^{\mathrm{op}}}(\Res Q,B).
\]
A projective bimodule resolution of \(B\) restricts to a
right-projective resolution of the free right module \(B\).
This proves the first vanishing.
For the second, tensor a left-projective resolution \(Q_\bullet\)
of \(DB\) over \(k\) with the right regular module \(B\).
It is a projective bimodule resolution of \(DB\otimes B\), and
\[
 \Hom_{B^e}(Q_\bullet\otimes B,DB)
 \cong\Hom_B(Q_\bullet,DB).
\]
Its positive cohomology vanishes since \({}_BDB\) is injective.
Both assertions pass to finite sums and summands.

Finally choose a projective resolution \(Q_B\to B\) and an
injective coresolution \(DB\to I_{DB}\), regarded as homological
complexes.  There are natural quasi-isomorphisms of Hom complexes
\begin{equation}\label{b:Hom-comparisons}
 \begin{split}
 \Hom^\bullet_{B^e}(B,P')&\longrightarrow
       \Hom^\bullet_{B^e}(Q_B,P'),\\
 \Hom^\bullet_{B^e}(P',DB)&\longrightarrow
       \Hom^\bullet_{B^e}(P',I_{DB}).
 \end{split}
\end{equation}
To verify them, filter by the finitely many terms of \(P'\).
For each successive quotient the claim is one of
\eqref{b:two-orthogonalities}, together with equality in degree
zero and vanishing of negative Ext between modules.  A shift of
a term merely shifts this comparison.  Hence terms in negative
degrees cause no exception, and induction on this finite
filtration proves \eqref{b:Hom-comparisons}.

Taking degree-zero cohomology shows that both derived arrows
have actual chain-map representatives
\(\alpha:B[0]\to P'\) and \(\beta:P'\to DB[0]\).
Their composite is the ordinary map \(\beta_0\alpha_0\), so it
factors through \(P'_0\in\add(DB\otimes B)\).
It equals \(v\) as an ordinary bimodule map, because
\(\Hom_{\mathrm D(B^e)}(B,DB)=\Hom_{B^e}(B,DB)\).
\end{proof}

The scalar trace calculation below is related to the trace and
corner-dimension calculation in
\cite[Example~3.2(1) and proof of Theorem~3.8]{LiuZhouZimmermann2012}.
Here the preceding lemma supplies the reduction from perfect
factorizations, and the calculation uses the split corner table
\eqref{a:corner-parity}.

\begin{lemma}\label{b:trace-obstruction}
The bimodule map
\begin{equation}\label{b:character-map}
 \eta:B\longrightarrow DB,\qquad
 \eta(a)=\chi_B(a)\chi_B,
\end{equation}
does not factor in the derived category through a perfect
\(B^e\)-complex.
\end{lemma}

\begin{proof}
The character identities make \(\eta\) a bimodule map.
By Lemma~\ref{b:ordinary-factorization}, a perfect factorization
would give an ordinary factorization through
\(\add(DB\otimes B)\).  We show that every such factorization
\(B\to DB\) sends \(1\) to a functional vanishing on \(p_f\).

First consider maps
\[
 B\xrightarrow{\alpha}DB\otimes B\xrightarrow{\beta}DB.
\]
Identify \(DB\otimes B\) with \(\Hom_k(B,B)\) by
\(\lambda\otimes c\mapsto[x\mapsto\lambda(x)c]\).
The element \(\alpha(1)\) satisfies
\(a\alpha(1)=\alpha(1)a\).
Under this identification, its corresponding linear map \(f_0\)
therefore satisfies
\[
 f_0(xa)=f_0(x)a\quad(x,a\in B).
\]
Thus \(f_0(x)=d_0x\), where \(d_0=f_0(1)\).
If \((c_i)\) is any vector-space basis of \(B\), with dual basis
\((c_i^*)\), then
\[
 \alpha(1)=\sum_i c_i^*\otimes d_0c_i.
\]
Right linearity of \(\beta\) reduces it to the left-module
endomorphism \(\lambda\mapsto\beta(\lambda\otimes1)\) of \(DB\).
Duality identifies these endomorphisms with endomorphisms of the
right regular module \(B\), which are left multiplications.
Thus the endomorphism of \(DB\) is right multiplication by some
\(b'\in B\).  Consequently
\[
 \beta(\lambda\otimes c)=\lambda b'c,\qquad
 (\beta\alpha(1))(p_f)
   =\sum_i c_i^*(b'd_0c_ip_f)
   =\tr_k\bigl(L_{b'd_0}R_{p_f}\bigr).
\]
Here \(L_a\) and \(R_a\) are left and right multiplication on
the vector space \(B\).  The order \(b'd_0\) follows from the
dual action \((\lambda b'c)(p_f)=\lambda(b'cp_f)\).

Write
\(b'd_0=r+\sum_v a_vp_v\), where \(r\in\rad B\) and the
\(p_v\) are its four vertex idempotents.
Left and right multiplications commute, and \(p_f^2=p_f\), so
\[
 (L_rR_{p_f})^n=L_{r^n}R_{p_f}=0
\]
for sufficiently large \(n\).
The radical part has trace zero.
For each \(v\), the operator \(L_{p_v}R_{p_f}\) is a projection
onto \(p_vBp_f\).  The corner dimensions of \(B\), obtained from
the corner table \eqref{a:corner-parity}, are even; in particular
\(\dim_k(p_vBp_f)=4\) for every \(v\).
Each projection therefore has trace zero in characteristic two.
Linearity gives \((\beta\alpha(1))(p_f)=0\).

A factorization through a finite sum of \(DB\otimes B\) is a sum
of such composites.  A factorization through a summand extends
to that finite sum by its split inclusion and projection.
The same vanishing thus holds throughout \(\add(DB\otimes B)\).
But \(\eta(1)(p_f)=\chi_B(p_f)=1\), a contradiction.
\end{proof}

\subsection{Lifts with nonzero evaluation}

We now apply the trace obstruction to obtain the maps required by
Proposition~\ref{b:bimodule}.

Put \(S=\Hom_k(X,X)\), with bimodule action
\[
 (atb)(x)=a\,t(bx).
\]
The representation gives a bimodule map \(\pi:E\to S\).
Since \(X\) is one-dimensional, \(S=k\) with both actions given
by \(\chi\).

\begin{lemma}\label{b:stable-adjunction}
For every finite side-projective bimodule \(L\), ordinary
tensor--Hom adjunction induces a natural isomorphism
\[
 \sHom_{E^e}(L,S)\ \cong\ \sHom_E(LX,X).
\]
This assertion does not require \(S\) to be side-projective.
\end{lemma}

\begin{proof}
The ordinary adjoint of \(f:L\to S\) is
\(\widehat f(l\otimes x)=f(l)(x)\).
A factorization of \(f\) through a finite \(E^e\)-projective
evaluates to a factorization through an \(E\)-projective.

Conversely embed \(L\) in a finite projective-injective
\(E^e\)-module \(J\).
The embedding splits as a right \(E\)-module map because
\(L_E\) is projective, hence injective over symmetric \(E\).
Thus \(LX\to JX\) is injective.
Suppose \(\widehat f\) factors as \(LX\to P\to X\) with \(P\)
finite \(E\)-projective.  Injectivity of \(P\) extends its first
factor to \(JX\to P\).  Composing with \(P\to X\) and applying
ordinary adjunction gives \(J\to S\) whose restriction to \(L\)
is \(f\).  This is a bimodule projective factorization.
The two factorization ideals therefore correspond exactly.
Naturality is inherited from ordinary adjunction.
\end{proof}

\begin{lemma}\label{b:evaluated-lift}
For every \(H\in k^\times\), there is a stable bimodule map
\(U_H[1]\to\cI\) with nonzero evaluated stable map
\(X[1]\to\cI X\).
\end{lemma}

\begin{proof}
Write \(U=U_H\) and
\(\phi=\pi\overline m:\cI\to S\).
All shifts of \(U\) and \(\cI\) have finite side-projective
representatives by Lemma~\ref{h:evaluation}.
Consider the finite-dimensional linear map
\[
 \phi_*:\sHom_{E^e}(U[1],\cI)
 \longrightarrow\sHom_{E^e}(U[1],S).
\]
Symmetric duality, Lemma~\ref{h:duality} applied over \(E^e\),
identifies its transpose with
\begin{equation}\label{b:transpose}
 \sHom_{E^e}(S,U)\longrightarrow\sHom_{E^e}(\cI,U),
 \qquad b\longmapsto b\phi.
\end{equation}
Indeed the fixed first argument is \(U[1]\), and its duality
shift is \(U[1][-1]=U\); naturality in the second argument gives
the displayed precomposition.

Let \(\zeta=f^*\otimes f^*\in E\), as in
Lemma~\ref{e:socle}.  Both its left and right actions
are the character \(\chi\).  The map
\[
 b_0:S\longrightarrow U,\qquad b_0(1)=\zeta
\]
is a bimodule map even for the twisted right action, because
\[
 \zeta\cdot b=\zeta\sigma_H(b)
       =\chi(\sigma_H(b))\zeta=\chi(b)\zeta.
\]
We will prove that \(b_0\phi\ne0\) stably.

Suppose otherwise.  The derived map \(b_0\pi m:K\to U\)
then has zero stable image and factors through a perfect
\(E^e\)-complex by Lemma~\ref{h:perfect-factor}.
To apply that lemma with projective models, one may take a
bounded-below projective resolution \(Q\to U\) with finite terms and lift the
global chain map \(K\to U\) to \(K\to Q\); its stable image is
the same map \(b_0\phi\).

Restriction of this factorization to \(B^e\) still has a perfect
intermediate object.  Indeed an \(E^e\)-free module restricts to
the outer bimodule \(E\otimes E\).  Tensoring the left and
right \(B\)-projective resolutions of \(E\) from
Proposition~\ref{a:dual-dimensions} gives a \(B^e\)-projective
resolution of length at most \(8\).  Finite sums and summands
have the same bound.  A bounded complex of these restricted
modules is perfect by its finite term filtration.

Let \(\operatorname{Bar}(B)\to B\) denote the bar resolution.
Inclusion of both outer \(B\)-factors in \(E\) defines a
\(B^e\)-chain map
\[
 \iota:\operatorname{Bar}(B)\longrightarrow\Res_{B^e}K,
\]
and hence a derived map \(B\to\Res_{B^e}K\).
Since \(\sigma_H\) fixes \(B\), the decomposition
\[
 \Res_{B^e}U
   =B\oplus(DC\otimes C)\oplus(C\otimes DC)\oplus DB
\]
is a decomposition of \(B\)-bimodules.  Let \(\rho\) be projection
to the last summand.  The composite
\[
 \operatorname{Bar}(B)\xrightarrow{\iota}\Res K
 \xrightarrow{\,b_0\pi m\,}\Res U\xrightarrow{\rho}DB
\]
is exactly the character map \(\eta\) of
\eqref{b:character-map} composed with the bar augmentation:
in degree zero it sends \(a\otimes c\) to
\(\chi_B(ac)\chi_B\), since \(\rho(\zeta)=\chi_B\), and in
higher degrees it is zero.
Thus our supposed factorization would make \(\eta\) factor
through a perfect \(B^e\)-complex.
Lemma~\ref{b:trace-obstruction} excludes this possibility.
We have proved \(b_0\phi\ne0\).

It follows from \eqref{b:transpose} that \(\phi_*\) is nonzero.
Choose \(\beta:U[1]\to\cI\) with \(\phi\beta\ne0\).
Lemma~\ref{b:stable-adjunction}, applied only to the
side-projective source \(U[1]\), identifies this nonzero class
with the composite
\[
 U[1]X\xrightarrow{\beta X}\cI X
       \xrightarrow{\overline m X}EX\cong X .
\]
This identity follows from ordinary adjunction and from the
evaluation of the representation \(\pi\).
The evaluated maps involving \(U[1],\cI,E\) are defined stably
by Lemma~\ref{h:evaluation}; no stable evaluation functor on
the nonside-projective module \(S\) is needed.
The composite is nonzero, so \(\beta X\) itself is nonzero.
The canonical shift comparison identifies \(U[1]X\) with
\((UX)[1]=X[1]\), proving the assertion.
\end{proof}

\begin{proof}[Proof of Proposition~\ref{b:bimodule}]
Take the fixed finite side-projective \(Y\) constructed in
Lemma~\ref{b:bar-model}.
Lemma~\ref{b:profile} proves its required Hom profile.
For any \(H\in k^\times\), shift the map of
Lemma~\ref{b:evaluated-lift} by \([-1]\) and identify
\(\cI[-1]\) with \(Y\) to obtain a stable map \(U_H\to Y\).
By the shift and evaluation compatibilities of
Lemma~\ref{h:evaluation}, its evaluation is a nonzero class
\(X\to YX\).
A stable map between the fixed finite modules is represented
by an ordinary bimodule map; choose one and call it \(g_H\).
All Hom spaces, dualities, and selections used above are over
the original field \(k\), so this requires no scalar extension.
The construction of \(Y\) did not depend on \(H\).
\end{proof}
